\documentclass[a4paper,twoside]{article}

% morewrites virtualises \newwrite so a document that combines many glossary
% groups (each costs one \write), biblatex (\bcfout), hyperref (\@outlinefile),
% the toc, and the LOT/LOF can exceed XeLaTeX's 16-register limit. Must be
% loaded before any package that allocates writes; that is why it sits ahead
% of xcolor here.
\usepackage{morewrites}

\usepackage[svgnames,dvipsnames,x11names]{xcolor}
\usepackage{graphicx}
% Babel must be loaded before any package that hooks into its localisation (notably
% the glossaries fragment in ``extra_packages``), otherwise the default acronym
% table title falls back to English regardless of ``language: french`` etc.
\usepackage[english]{babel}
% Hyperref must be loaded before ``glossaries`` (pulled in via ``extra_packages``)
% so glossaries can wrap glossary page numbers in ``\hyperpage`` for clickable
% backreferences. Loading it after glossaries leaves the page numbers as plain
% text without hyperlinks.
\usepackage[colorlinks=true, linkcolor=NavyBlue, urlcolor=NavyBlue, citecolor=NavyBlue]{hyperref}
\usepackage[a4paper]{geometry}
\usepackage{ts-fonts}
\usepackage{marginnote}
\usepackage{amsmath}
\usepackage{float}
\usepackage{seqsplit}
% Default margin-note styling. Notes are typeset at ``\footnotesize`` so they
% fit narrow margins, and ``\sloppy`` + a generous ``\emergencystretch`` plus
% lowered hyphenation penalties let the line breaker stretch / hyphenate
% rather than overshoot the box: a long technical word (URL, German compound,
% French ``anticonstitutionnellement``…) hyphenates instead of poking past
% the page edge. Override with ``\renewcommand*{\marginfont}{…}`` in a
% custom preamble snippet to disable any of these.
\renewcommand*{\marginfont}{%
  \footnotesize
  \sloppy
  \emergencystretch=1em
  \hyphenpenalty=50
  \exhyphenpenalty=50
}
% Clamp ``\marginparwidth`` so margin notes never overflow the page edge,
% even when the user's geometry reserves less horizontal space than LaTeX's
% default ``\marginparwidth`` (~3.9 cm). We compute the space actually
% available on each side (recto / verso) at ``\AtBeginDocument`` time —
% once ``geometry`` has finalised ``\textwidth`` / ``\oddsidemargin`` /
% ``\evensidemargin`` — and keep the smaller of the two so the note fits on
% both sides in twoside documents. A 6 pt safety buffer is subtracted on top
% of ``\marginparsep`` so a slightly-too-long internal line still has room
% before clipping the page edge.
\newlength{\tsmarginparavail}
\newlength{\tsmarginparalt}
\newlength{\tsmarginparbuf}
\setlength{\tsmarginparbuf}{6mm}
\makeatletter
\AtBeginDocument{%
  \setlength{\tsmarginparavail}{\dimexpr\paperwidth-\textwidth-1in-\oddsidemargin-\marginparsep-\tsmarginparbuf\relax}%
  \ifdim\tsmarginparavail<\z@\setlength{\tsmarginparavail}{\z@}\fi
  \if@twoside
    \setlength{\tsmarginparalt}{\dimexpr1in+\evensidemargin-\marginparsep-\tsmarginparbuf\relax}%
    \ifdim\tsmarginparalt<\z@\setlength{\tsmarginparalt}{\z@}\fi
    \ifdim\tsmarginparalt<\tsmarginparavail
      \setlength{\tsmarginparavail}{\tsmarginparalt}%
    \fi
  \fi
  \ifdim\tsmarginparavail<\marginparwidth
    \setlength{\marginparwidth}{\tsmarginparavail}%
  \fi
}
\makeatother
\usepackage{ts-code}

\usepackage{lastpage} % pour le label LastPage
\usepackage{refcount} % pour \getpagerefnumber
\makeatletter
\def\ps@plain{%
  \let\@oddhead\@empty
  \let\@evenhead\@empty
  \def\@oddfoot{%
    \hfil
    % Afficher le numéro seulement si le document a plus d'une page
    \ifnum\getpagerefnumber{LastPage}>1\relax
      \thepage
    \fi
    \hfil}%
  \let\@evenfoot\@oddfoot
}
\makeatother

\title{Margin Notes\\[0.5em]\large The \texttt{\{margin\}[…]} inline extension}
\author{TeXSmith}\date{}

\begin{document}

\maketitle

\section{Margin Notes}\label{margin-notes}

TeXSmith's margin-note extension adds a single inline shorthand to the
unified \texttt{\{keyword\}[content]} family already used by \texttt{\{index\}[term]} and
\texttt{\{latex\}[payload]}. The syntax is:

\begin{code}{text}{}{}
\begin{Verbatim}[commandchars=\\\{\},breaklines, breakanywhere, commandchars=\\\{\}]
\PYZob{}margin\PYZcb{}[note text]\PYZob{}side?\PYZcb{}
\end{Verbatim}
\end{code}
where \texttt{side} is an optional single-letter suffix --- \texttt{l}, \texttt{r}, \texttt{o}, or \texttt{i} ---
selecting a margin explicitly. It compiles down to \texttt{\textbackslash{}marginnote\{…\}} from the
\texttt{marginnote} \LaTeX{} package (auto-loaded on first use).

\subsection{Default placement}\label{default-placement}

No suffix means the document's default side. In a \texttt{oneside} layout that's
the right margin; in a \texttt{twoside} layout, the outer margin (so the note
flips automatically between recto and verso pages).

Most readers appreciate a short quip\marginnote{default note} that sits
beside the paragraph without interrupting the flow. When margin notes are
kept brief and self-contained, they feel like a friendly aside rather than
a distraction.

\subsection{Forced side}\label{forced-side}

Force the left margin with \texttt{\{l\}} and the right margin with \texttt{\{r\}}. Because
the switch is scoped to a \LaTeX{} group, subsequent unqualified notes still
follow the document's default placement.

Some diagrams read better when labelled on the left{\reversemarginpar\marginnote{left-hand
pointer}}, while running commentary fits the right-hand
margin\marginnote{right-hand commentary} more naturally. Mixing the two in
close succession is fine --- each note is independent.

\subsection{Inline formatting}\label{inline-formatting}

Margin notes pass through the full inline-Markdown parser, so they support
the usual \textbf{bold}, \emph{italic}, \texttt{inline code}, and \href{https://example.org}{links}.

Viscoelastic materials can be approximated as linear\marginnote{\textbf{Hooke's}
law applies at small strain where the stress is \emph{proportional} to the
strain} under moderate stress, but non-linear effects dominate above the
yield point.

\subsection{Longer notes}\label{longer-notes}

The package handles multi-sentence notes gracefully. They line-wrap in the
margin and do not disturb the main text's baseline
grid.\marginnote{Longer notes wrap over several lines and keep flowing down
the margin. Keep them concise so the reader can scan them without losing
the thread of the main text.}

Below, a denser example mixing sides and inline styles:

\begin{itemize}
\item{} Classical mechanics{\reversemarginpar\marginnote{see Newton, \emph{Principia}, 1687}} predates
  the calculus of variations.
\item{} Modern formulations rely on Lagrangians\marginnote{or Hamiltonians for
  energy-based analyses} and variational principles.
\item{} Quantum mechanics\marginnote{introduces non-commuting
  observables} follows in the 20\textsuperscript{th} century.

\end{itemize}
\subsection{Inner / Outer}\label{inner-outer}

In a \texttt{twoside} layout --- like this article --- \texttt{\{o\}} (outer) and \texttt{\{i\}} (inner)
are MVP aliases of \texttt{\{r\}} and \texttt{\{l\}} respectively. On single-side documents
they behave identically to their counterparts.

\end{document}
